\section{Entropy clipping and large representation dimensions}
\label{l:replica-section}

Small total variation error is one way to obtain a useful coset cap.
A weaker entropy estimate can also suffice when only representations
above a prescribed dimension threshold need to be controlled.  The
cutoff below preserves the original, unnormalized retained measure.
Its loss of mass is separated from its Fourier estimate.

For probabilities $p,u$ on a finite set, with $u$ positive, write
\[
 D(p\Vert u)=\sum_xp(x)\log\frac{p(x)}{u(x)};
\]
zero summands have value zero and logarithms are natural.
For $G=S_n$ and $H\le G$, write $w_H(gH)$ for the mass of the
left coset $gH$ under a probability $w$ on $G$, and $U_{G/H}$ for
the uniform law on left cosets.  These are the cosets that record the
images of a list whose pointwise stabilizer is $H$.

\begin{theorem}[Entropy clipping and isotypic transfer]
\label{l:replica-clipping}
Partition $\{1,\ldots,n\}$ into $q\ge1$ nonempty blocks, and let
$H_i$ be the subgroup supported on block $i$.  Suppose a probability
$w$ on $S_n$ satisfies
\begin{equation}\label{l:replica-entropy-assumption}
 \sum_{i=1}^qD(w_{H_i}\Vert U_{G/H_i})\le K.
\end{equation}
For every $t>0$ there is a subprobability $v\le w$, of mass $m$,
such that
\begin{align}
 1-m&\le\frac{K+q/e}{t},\label{l:replica-mass}\\
 v(gH_i)&\le e^t\frac{|H_i|}{|G|}\quad(g\in G,\ 1\le i\le q),
 \label{l:replica-cap}\\
 \sum_{h\in H_i}(\check v*v)(h)
 &\le e^t\frac{|H_i|}{|G|},\label{l:replica-collision}\\
 \|\widehat v(\lambda)\|_{\mathrm{HS}}^2
 &\le qC_u e^tD_\lambda^{-1+(u+2)/q}
        \quad(\lambda\vdash n, u>0).
 \label{l:replica-hs}
\end{align}
Here $\check v(g)=v(g^{-1})$, $C_u$ is the constant in
Theorem~\ref{l:degree-all-powers}, the Fourier transform is the mass
transform $\widehat v(\lambda)=\sum_gv(g)\rho_\lambda(g)$, and
the Hilbert--Schmidt norm is unnormalized.
\end{theorem}
\begin{proof}
For a density $p$ relative to a probability $u$, the negative part of
the logarithm has expectation at most $1/e$ under $pu$, since
$p\max(-\log p,0)\le1/e$.  Consequently
\[
 \mathbb E_{pu}(\log p)_+\le D(pu\Vert u)+1/e.
\]
For each $i$ let $p_i$ be the density of $w_{H_i}$ relative to
$U_{G/H_i}$.  Restrict $w$ to those $g$ for which
$p_i(gH_i)\le e^t$ for every $i$, and call the result $v$.
Markov's inequality and a union bound prove \eqref{l:replica-mass};
the restriction proves \eqref{l:replica-cap}.  Zero-density cosets
have zero $w$-mass and create no logarithmic exception.

The convolution $\check v*v$ is the law, with its subprobability
mass, of $g_1^{-1}g_2$.  Membership in $H_i$ means that $g_1,g_2$
belong to the same left coset.  Thus
\[
 \sum_{h\in H_i}(\check v*v)(h)
 =\sum_{gH_i}v(gH_i)^2
 \le e^t\frac{|H_i|}{|G|}\sum_{gH_i}v(gH_i)
 \le e^t\frac{|H_i|}{|G|}.
\]

We spell out the isotypic step to show exactly what the collision cap
controls.  Put $E_\lambda=\widehat v(\lambda)^*\widehat v(\lambda)$.
For an irreducible $\tau$ of $H_i$ of dimension $d_\tau$ and
character $\chi_\tau$, let $P_{\lambda;\tau}^{H_i}$ be the full
$\tau$-isotypic projection in $\rho_\lambda|_{H_i}$, including all
multiplicity copies.  Fourier inversion and the central projection
formula give
\begin{equation}\label{l:replica-isotypic-identity}
 \sum_{h\in H_i}(\check v*v)(h)\chi_\tau(h)
 =\frac{|H_i|}{|G|d_\tau}
   \sum_{\lambda\vdash n}D_\lambda
         \operatorname{tr}(E_\lambda P_{\lambda;\tau}^{H_i}).
\end{equation}
Indeed $\sum_{h\in H_i}\chi_\tau(h)\rho_\lambda(h^{-1})
=(|H_i|/d_\tau)P_{\lambda;\tau}^{H_i}$.
Every trace on the right is nonnegative, since $E_\lambda$ and
$P_{\lambda;\tau}^{H_i}$ are positive.  Since
$|\chi_\tau(h)|\le d_\tau$, \eqref{l:replica-collision} yields
\begin{equation}\label{l:replica-one-isotype}
 \operatorname{tr}(E_\lambda P_{\lambda;\tau}^{H_i})
       \le e^t\frac{d_\tau^2}{D_\lambda}.
\end{equation}
This is the collision form of Lemma~\ref{l:central-isotypic}.

The block subgroups form a direct product.  Every irreducible
constituent of its restriction has factor dimensions whose product
is at most $D_\lambda$.  At least one factor dimension is therefore
at most $D_\lambda^{1/q}$.  The commuting isotypic projections satisfy
\[
 I\le\sum_{i=1}^q\ \sum_{\tau:d_\tau\le D_\lambda^{1/q}}
                      P_{\lambda;\tau}^{H_i}.
\]
There are at most $C_uD_\lambda^{u/q}$ such types in each subgroup:
each contributes at least $D_\lambda^{-u/q}$ to its inverse-$u$
degree sum.  Taking traces against $E_\lambda$ and using
\eqref{l:replica-one-isotype} proves \eqref{l:replica-hs}.
\end{proof}

\begin{corollary}[Selecting complementary lists by averaging]
\label{l:replica-partition-selection}
Let $q$ divide $n$, put $r=n/q$, and suppose the image law $w_I$ of
a uniformly chosen ordered list $I$ of $n-r$ distinct inputs satisfies
\[
 \mathbb E_I D(w_I\Vert U_I)\le K_0,
\]
where $U_I$ is uniform on ordered distinct output lists of that length.
There is a partition into $q$ blocks of size $r$ for which
\eqref{l:replica-entropy-assumption} holds with $K=qK_0$.
\end{corollary}
\begin{proof}
Choose the equal partition uniformly.  Each block complement is
uniform among sets of size $n-r$.  Entropy of an image list is
unchanged by reordering its coordinates.  Therefore the expected sum
of the $q$ complement entropies is at most $qK_0$; one deterministic
partition has sum no greater than that expectation.  Agreement on
the complement is exactly membership in the same left coset of its
block subgroup.
\end{proof}

\begin{corollary}[The logarithmic cutoff with 256 blocks]
\label{l:replica-large-dimension}
Let $n=2^d$ and suppose, for each sufficiently large integer $d$, that
a probability $w$ has an equal partition into $q=256$ blocks with
total complement entropy $O(d^4)$.  There is $v\le w$ of mass
$1-O(d^{-2})$ such that
\begin{equation}\label{l:replica-34-exponent}
 \|\widehat v(\lambda)\|_{\mathrm{HS}}^2
 \le256 C_{32}e^{d^6}D_\lambda^{-1+34/256}.
\end{equation}
In particular, whenever $\log D_\lambda>d^8$, this is at most
$D_\lambda^{-1/2}$ for all sufficiently large $d$, uniformly in
$\lambda$.
\end{corollary}
\begin{proof}
Apply Theorem~\ref{l:replica-clipping} with $t=d^6$ and $u=32$.
The loss bound is $O(d^4)/d^6=O(d^{-2})$.
The logarithm of the prefactor is $d^6+\log(256C_{32})$, while
\[
 \left(\frac12-\frac{34}{256}\right)\log D_\lambda
 >\left(\frac12-\frac{34}{256}\right)d^8.
\]
The latter eventually exceeds the former, proving the final bound.
\end{proof}

The hypotheses here concern only the entropy of a group law.
Controlling dimensions below the cutoff requires an additional input
on that law; the clipping theorem itself makes no assertion about
those representations.
